\fichatitulo{33 · ESCALA}{Preprint v1 · 2020}{Scaling Laws for Neural Language Models}
{\small Kaplan et al. \textit{Scaling Laws for Neural Language Models}. Preprint v1 · 2020. \href{https://arxiv.org/pdf/2001.08361v1}{Fonte consultada}.\par}

\campo{Problema e objetivo}
Como a perda de modelagem varia com parâmetros, dados e computação?

\campo{Principal contribuição · síntese interpretativa}
Descreve relações empíricas de escala para a perda de entropia cruzada de modelos de linguagem.

\campo{Método / natureza da evidência}
Ajuste de tendências observadas ao variar tamanho do modelo, quantidade de dados e orçamento computacional.

\campo{Resultado ou argumento principal}
Discute alocação de computação que favorece modelos maiores e interrupção do treinamento antes da convergência nas condições analisadas.

\campo{Excertos-chave · seleção literal}
\begin{itemize}
\excerto{1}{cross-entropy loss}{Resumo.}
\excerto{2}{stopping significantly before convergence}{Resumo.}
\end{itemize}

\campo{Limitações e análise crítica}
Análise da ficha: tendências empíricas dependem do regime experimental. Reduzir perda não equivale a melhorar fidelidade de RAG; confrontar pressupostos com Hoffmann et al.

\campo{Aproveitamento na dissertação · proposta}
Fundamentação: contextualizar custo e escala sem usar número de parâmetros como substituto da avaliação de tarefa.

\registro{Fonte consultada nesta preparação: Resumo e introdução. Chave: \texttt{kaplanEtAl2020scalingLaws}.}
